\subsection{EXACT SEQUENCES}

DEFINITION 5. Let F, G, H be three A-modules; let f be a homomorphism of F into G and g a homomorphism of G into H. The ordered pair \((f, g)\) is called an exact sequence if

\[
\bar {g} ^ {- 1} (0) = f (\mathbf {F})
\]

in other words, if the kernel of g is equal to the image of f.

The diagram

\[
\mathrm{F} \xrightarrow {f} \mathrm{G} \xrightarrow {g} \mathrm{H}\tag{12}
\]

is also called an exact sequence.

We consider similarly a diagram consisting of four A-modules and three homomorphisms:

\[
\mathrm{E} \xrightarrow {f} \mathrm{F} \xrightarrow {g} \mathrm{G} \xrightarrow {h} \mathrm{H}.\tag{13}
\]

This diagram is called exact at F if the diagram \(E \xrightarrow{f} F \xrightarrow{g} G\) is exact; it is called exact at G if \(F \xrightarrow{g} G \xrightarrow{h} H\) is exact. If diagram (13) is exact at F and at G, it is simply called exact, or an exact sequence. Exact sequences with an arbitrary number of terms are defined similarly.

Remark. (1) If the ordered pair \((f, g)\) is an exact sequence, then \(g \circ f = 0\) ; but of course this property does not characterize exact sequences, for it only means that the image of f is contained in the kernel of g.

In the statements below, E, F, G denote A-modules, 0 the A-module reduced to its identity element; the arrows represent A-module homomorphisms. As there is only one homomorphism from the module 0 to a module E (resp. of E to 0), there is no point in giving these homomorphisms a name in the exact sequences where they appear.

PROPOSITION 3. (a) For

\[
0 \longrightarrow \mathbf {E} \stackrel {f} {\longrightarrow} \mathbf {F}
\]

to be an exact sequence, it is necessary and sufficient that f be injective.

\begin{enumerate}
\def\labelenumi{(\alph{enumi})}
\setcounter{enumi}{1}
\tightlist
\item
  For
\end{enumerate}

\[
\mathbf {E} \xrightarrow {f} \mathbf {F} \longrightarrow 0
\]

to be an exact sequence, it is necessary and sufficient that \(f\) be surjective.

\begin{enumerate}
\def\labelenumi{(\alph{enumi})}
\setcounter{enumi}{2}
\tightlist
\item
  For
\end{enumerate}

\[
0 \longrightarrow \mathrm{E} \stackrel {f} {\longrightarrow} \mathrm{F} \longrightarrow 0
\]

to be an exact sequence, it is necessary and sufficient that \(f\) be bijective (in other words (no. 2, Proposition 2) that \(f\) be an isomorphism of \(\mathbf{E}\) onto \(\mathbf{F}\) ).

\begin{enumerate}
\def\labelenumi{(\alph{enumi})}
\setcounter{enumi}{3}
\tightlist
\item
  If \(\mathbf{F}\) is a submodule of \(\mathbf{E}\) and \(i: \mathbf{F} \to \mathbf{E}\) is the canonical injection, \(p: \mathbf{E} \to \mathbf{E} / \mathbf{F}\) the canonical homomorphism, the diagram
\end{enumerate}

\[
0 \longrightarrow \mathrm{F} \stackrel {i} {\longrightarrow} \mathrm{E} \stackrel {p} {\longrightarrow} \mathrm{E} / \mathrm{F} \longrightarrow 0\tag{14}
\]

is an exact sequence.

\begin{enumerate}
\def\labelenumi{(\alph{enumi})}
\setcounter{enumi}{4}
\tightlist
\item
  If \(f: \mathbf{E} \to \mathbf{F}\) is a homomorphism, the diagram
\end{enumerate}

\[
0 \longrightarrow f ^ {- 1} (0) \stackrel {i} {\longrightarrow} \mathrm{E} \stackrel {f} {\longrightarrow} \mathrm{F} \stackrel {p} {\longrightarrow} \mathrm{F} / f (\mathrm{E}) \rightarrow 0
\]

(where i is the canonical injection and p the canonical surjection) is an exact sequence.

The proposition follows immediately from the definitions and Proposition 2 of no. 2.

Remarks. (2) To say that there is an exact sequence

\[
0 \longrightarrow \mathrm{E} \stackrel {f} {\longrightarrow} \mathrm{F} \stackrel {g} {\longrightarrow} \mathrm{G} \longrightarrow 0
\]

means that \(f\) is injective, \(g\) surjective and that the canonical bijection associated with \(g\) is an isomorphism of \(\mathrm{F} / f(\mathbf{E})\) onto \(\mathbf{G}\) . It is also said that the triple \((\mathbf{F}, f, g)\) is an extension of the module \(\mathbf{G}\) by the module \(\mathbf{E}\) (I, § 6, no. 7).

\begin{enumerate}
\def\labelenumi{(\arabic{enumi})}
\setcounter{enumi}{2}
\tightlist
\item
  If there is an exact sequence with 4 terms
\end{enumerate}

\[
\mathrm{E} \xrightarrow {f} \mathrm{F} \xrightarrow {g} \mathrm{G} \xrightarrow {h} \mathrm{H}
\]

the cokernel of \(f\) is \(\mathrm{F} / f(\mathrm{E}) = \mathrm{F} / g^{-1}(0)\) and the kernel of \(h\) is \(g(\mathrm{F})\) ; the canonical bijection associated with \(g\) is therefore an isomorphism

\[
\operatorname{Coker} f \cong \operatorname{Ker} h.
\]

\begin{enumerate}
\def\labelenumi{(\arabic{enumi})}
\setcounter{enumi}{3}
\tightlist
\item
  Consider an ordered pair of A-module homomorphisms
\end{enumerate}

\[
\mathrm{E} \xrightarrow {f} \mathrm{F} \xrightarrow {g} \mathrm{G}.\tag{15}
\]

For diagram (15) to be an exact sequence, it is necessary and sufficient that there exist two A-modules S, T and homomorphisms \(a:E\to S\) , \(b:S\to F\) , \(c:F\to T\) , \(d:T\to G\) such that the three sequences

\[
\left\{ \begin{array}{c} \mathrm{E} \xrightarrow {a} \mathrm{S} \longrightarrow 0 \\ 0 \longrightarrow \mathrm{S} \xrightarrow {b} \mathrm{F} \xrightarrow {c} \mathrm{T} \longrightarrow 0 \\ 0 \longrightarrow \mathrm{T} \xrightarrow {d} \mathrm{G} \end{array} \right.\tag{16}
\]

are exact and \(f = b \circ a\) and \(g = d \circ c\) .

If (15) is an exact sequence, then take \(\mathbf{S} = f(\mathbf{E}) = \overline{g}^{-1}(0)\) and \(\mathbf{T} = g(\mathbf{F})\) , \(b\) and \(d\) being the canonical injections and \(a\) (resp. \(c\) ) the homomorphism with the same graph as \(f\) (resp. \(g\) ). Conversely, if \(\mathbf{S}, \mathbf{T}, a, b, c, d\) satisfy the above conditions, then \(f(\mathbf{E}) = b(a(\mathbf{E})) = b(\mathbf{S})\) and \(g^{-1}(0) = c^{-1}(d^{-1}(0)) = c^{-1}(0)\) , hence the exactness of (16) shows that \(f(\mathbf{E}) = \overline{g}^{-1}(0)\) .

The use of explicit letters to denote homomorphisms in an exact sequence is often dispensed with when it is not necessary for the arguments.

Remark. (5) The definition of an exact sequence extends immediately to groups which are not necessarily commutative; in this case of course multiplicative notation is used with 0 replaced by 1 in the formulae (if no confusion arises). Parts (a), (b), (c) of Proposition 3 are still valid and also (d) when F is a normal subgroup of E. Remark 2 and Proposition 3(e) hold on condition that \(f(\mathrm{E})\) is a normal subgroup of F; Remarks 3 and 4 are valid without modification.

